\documentclass[12pt,a4paper]{amsart}

\usepackage{amsmath,amssymb,amsthm}
\usepackage{mathtools}
\usepackage{hyperref}
\usepackage{geometry}
\geometry{margin=2.5cm}

% -------------------------------------------------------
\newtheorem{theorem}{Theorem}[section]
\newtheorem{lemma}[theorem]{Lemma}
\newtheorem{corollary}[theorem]{Corollary}
\newtheorem{proposition}[theorem]{Proposition}
\theoremstyle{definition}
\newtheorem{definition}[theorem]{Definition}
\newtheorem{remark}[theorem]{Remark}

\newcommand{\Z}{\mathbb{Z}}

% -------------------------------------------------------
\begin{document}
% -------------------------------------------------------

\title{Every Giuga Pseudoprime is Squarefree:\\ A Short Proof}

\author{Kurt Ingwer}
\address{Specialist Mathematics Project, Version Build 44}
\email{specialist-maths@localhost}
\date{\today}

\begin{abstract}
A \emph{Giuga pseudoprime} is a composite integer satisfying a primality criterion
due to Giuga~(1950).  We give a short, self-contained proof that every Giuga
pseudoprime is necessarily squarefree.  The argument is elementary: if $p^2 \mid n$,
then the weak Giuga condition forces $p \mid -1$, which is impossible.  This result
is foundational for the study of Giuga pseudoprimes and justifies restricting
attention to squarefree composites from the outset.
\end{abstract}

\maketitle

% -------------------------------------------------------
\section{Introduction}
% -------------------------------------------------------

In 1950, Giuga \cite{Giuga1950} conjectured that a positive integer $n > 1$ is prime
if and only if $\sum_{k=1}^{n-1} k^{n-1} \equiv -1 \pmod{n}$.
A composite counterexample would be called a \emph{Giuga pseudoprime}.

Borwein, Borwein, Girgensohn, and Parnes \cite{Borwein1996} characterised such
numbers as composites satisfying, for every prime $p \mid n$:
\begin{enumerate}
  \item[\textup{(W)}] $p \mid \dfrac{n}{p} - 1$ \quad \emph{(weak condition)},
  \item[\textup{(S)}] $(p-1) \mid \dfrac{n}{p} - 1$ \quad \emph{(strong condition)}.
\end{enumerate}

The present note proves that any integer satisfying even the weak condition~(W) alone
must be squarefree.

% -------------------------------------------------------
\section{Main Result}
% -------------------------------------------------------

\begin{theorem}
\label{thm:sqfree}
Every Giuga pseudoprime is squarefree.  More precisely: if a composite integer~$n$
satisfies the weak Giuga condition~\textup{(W)} for all primes $p \mid n$, then
$n$ is squarefree.
\end{theorem}

\begin{proof}
Suppose for contradiction that $p^2 \mid n$ for some prime~$p$.
Then $p \mid \tfrac{n}{p}$, so
\[
  \frac{n}{p} \equiv 0 \pmod{p},
\]
which gives $\tfrac{n}{p} - 1 \equiv -1 \pmod{p}$.
The weak condition~(W) requires $p \mid \tfrac{n}{p} - 1$, i.e.\ $p \mid -1$.
But $p \ge 2$, so $p \nmid 1$ and hence $p \nmid -1$ — a contradiction.
Therefore no prime~$p$ can satisfy $p^2 \mid n$, i.e.\ $n$ is squarefree.
\end{proof}

\begin{remark}
The proof uses only condition~(W); the strong condition~(S) is not needed.
Hence the stronger statement holds: any composite satisfying only~(W) for
all prime factors is already squarefree.
\end{remark}

\begin{remark}
The squarefreeness is immediate from the structure of the weak condition.
Intuitively, if $p^2 \mid n$, then $\tfrac{n}{p}$ is divisible by~$p$, so
$\tfrac{n}{p} - 1$ is congruent to~$-1$ modulo~$p$, and no prime can
divide~$-1$.
\end{remark}

% -------------------------------------------------------
\section{Consequences}
% -------------------------------------------------------

\begin{corollary}
Every Giuga pseudoprime $n$ has the form $n = p_1 p_2 \cdots p_k$
with $p_1 < p_2 < \cdots < p_k$ distinct primes.
\end{corollary}

\begin{corollary}
For a Giuga pseudoprime $n = p_1 \cdots p_k$, the quotients
$\tfrac{n}{p_i}$ are products of distinct primes different from~$p_i$.
\end{corollary}

The squarefreeness established here is the starting point for further structural
results on Giuga pseudoprimes, including the proofs that no such number can have
fewer than four prime factors \cite{Ingwer2026a} and the related CRT obstructions
for Giuga--Carmichael numbers \cite{Ingwer2026c}.

% -------------------------------------------------------
\section{Numerical Illustration}
% -------------------------------------------------------

The known \emph{Giuga numbers} (satisfying only the weak condition) are
$\{30, 858, 1722, 66198, \ldots\}$.  All are squarefree:
\begin{align*}
  30 &= 2 \cdot 3 \cdot 5, \\
  858 &= 2 \cdot 3 \cdot 11 \cdot 13, \\
  1722 &= 2 \cdot 3 \cdot 7 \cdot 41, \\
  66198 &= 2 \cdot 3 \cdot 11 \cdot 17 \cdot 59.
\end{align*}
Theorem~\ref{thm:sqfree} guarantees this: any composite satisfying~(W) is squarefree.
One verifies directly that each of these numbers fails~(S) for at least one prime factor,
so none is a Giuga pseudoprime.

% -------------------------------------------------------
\begin{thebibliography}{10}

\bibitem{Giuga1950}
G.~Giuga,
\emph{Su una presumibile propriet\`a caratteristica dei numeri primi},
Ist.\ Lombardo Sci.\ Lett.\ Rend.\ Cl.\ Sci.\ Mat.\ Nat.\ \textbf{83}
(1950), 511--528.

\bibitem{Borwein1996}
J.~M.~Borwein, P.~B.~Borwein, R.~Girgensohn, and S.~Parnes,
\emph{Giuga's conjecture on primality},
Amer.\ Math.\ Monthly \textbf{103} (1996), no.\ 1, 40--50.

\bibitem{Ingwer2026a}
K.~Ingwer,
\emph{No composite number with exactly three prime factors is a Giuga pseudoprime},
Specialist Mathematics Project (2026).

\bibitem{Ingwer2026c}
K.~Ingwer,
\emph{On the incompatibility of Giuga pseudoprimes and Carmichael numbers},
Specialist Mathematics Project (2026).

\end{thebibliography}

\end{document}
